\begin{figure}[h!]
  \begin{subfigure}[c]{1.0\linewidth}
    \centering
    \includegraphics[height=.08\textheight]{persistent-naming-pb-1}
    \caption{}
    \label{fig:persistent-naming-pb-1}
  \end{subfigure}
  \begin{subfigure}[c]{1.0\linewidth}
    \centering
    \includegraphics[height=.08\textheight]{persistent-naming-pb-2}
    \caption{}
    \label{fig:persistent-naming-pb-2}
  \end{subfigure}
  \caption{%
    Spécification paramétrique. %
    (a) Évaluation. %
    (b) Réévaluation. %
  }
  \label{fig:persistent-naming-pb}
\end{figure}

Un problème bien connu en CAO, et plus particulièrement dans la modélisation
paramétrée, est celui de la persistance des références topologiques. %
En effet, à l'issue d'un processus de modélisation, une spécification
paramétrique ainsi que des références vers les entités topologiques impactés par
les opérations successivement appliquées sont créées. %
Ceci à pour objectif de rendre possible le rejeu d'un processus de
modélisation. %
Toutefois, lors du rejeu d'un processus de modélisation et lorsque la
spécification initiale a été éditée, la topologie est elle-même susceptible
d'être modifiée voire de ne plus exister et les références précédemment créées
ne sont alors plus valides.

La Figure~\ref{fig:persistent-naming-pb} illustre ce problème. %
Au jeu initial (Figure~\ref{fig:persistent-naming-pb-1}), la première opération
crée un cylindre et sa face du haut est nommée $f_{1}$. %
L'opération suivante creuse une rainure dans $f_{1}$ qui est alors scindée. %
Les deux faces issues de cette scission sont nommées $f_{2}$ et $f_{3}$ et
$f_{1}$ n'est plus référençable. %
Et enfin, un second cylindre est créé sur la face $f_{2}$. %
Au rejeu (Figure~\ref{fig:persistent-naming-pb-2}), le seul changement apporté
est la longueur de la rainure creusée dans $f_{1}$. %
Néanmoins, ce changement fait que $f_{1}$ n'est pas scindée au rejeu et $f_{2}$
n'existe pas. %
La référence à $f_{2}$ n'est donc plus valide et l'application de l'opération
créant le second cylindre ne produira pas le résultat attendu. %


Pour répondre à cette problématique, un système de nommage persistant des
entités topologiques référencées doit alors être mis en œuvre afin de garantir
que des résultats cohérents seront produits à l'issu d'un rejeu. %
Un tel système fonctionne en produisant, pour toute entité topologique
référencée, un nom unique et de les redécouvrir au rejeu par appariement. %
De tels systèmes ont déjà été présentés par \cite{kripac1995mechanism,
  CAPOYLEAS199617, wu2001face}. %
La déclaration de caractéristiques de faces (ou feature faces) est proposée par
Bidarra et al. dans \cite{bidarra2005feature}. %
Baba-Ali et al. \cite{baba2009edge} proposent un système basé sur l'appariement
d'arêtes pour résoudre l'ambigüité du nommage des arêtes en cas intersections
avec des faces non-planaires. %


% \begin{itemize}
%   \item kripac
%   \item capoyleas
%   \item wu
%   \item bidarra
%   \item baba-ali
%
%   \item marcheix
%   \item agbodan
%   \item farjana
%   \item cheon
%   \item cardot
% \end{itemize}


%%% Local Variables:
%%% mode: latex
%%% TeX-master: "../main"
%%% End:
